% exercise_set_id: set-derivative-rules-001
% mode: computation
% topic: derivatives
% source_image: https://github.com/wsollers/lra-proof-vault/tree/master/volume-vi/calculus/derivatives/exercises/derivative-computation-p-set-2026-06-04
% status: memorialized

\paragraph{Exercise Set: Derivative Computation}
\textbf{Stable set ID:} \texttt{set-derivative-rules-001}.
\textbf{Mode:} computation.
\textbf{Topic:} derivatives.
\textbf{Status:} memorialized.

\ProofVaultURL{https://github.com/wsollers/lra-proof-vault/tree/master/volume-vi/calculus/derivatives/exercises/derivative-computation-p-set-2026-06-04}

\begin{enumerate}
  \item \textbf{Stable item ID:} \texttt{ex-deriv-rules-001}.

  \textbf{Exercise.}
  Find the derivative of
  \[
    y = -7x^5 - 4x^3 + 9x - 2.
  \]

  \textbf{Solution.}
  By the power rule and linearity,
  \[
    \frac{dy}{dx}
      = -35x^4 - 12x^2 + 9.
  \]

  \textbf{Verification.}
  The result is corrected. The memorialized image has a negative leading
  coefficient, so the derivative keeps that sign.

  \item \textbf{Stable item ID:} \texttt{ex-deriv-rules-002}.

  \textbf{Exercise.}
  Find the derivative of
  \[
    f(x) = x^2 \sin x.
  \]

  \textbf{Solution.}
  By the product rule,
  \[
    f'(x) = 2x\sin x + x^2\cos x.
  \]

  \textbf{Verification.}
  The result is corrected. The problem was initially copied as a quotient, but
  the memorialized exercise is the product \(x^2\sin x\).

  \textbf{Rework note.}
  Preserve the correction in the ledger so later refactors do not recover the
  discarded quotient transcription.

  \item \textbf{Stable item ID:} \texttt{ex-deriv-rules-003}.

  \textbf{Exercise.}
  Find the derivative of
  \[
    y = \frac{x^2 + 1}{x - 3}.
  \]

  \textbf{Solution.}
  By the quotient rule,
  \[
    \frac{dy}{dx}
      =
      \frac{(x - 3)(2x) - (x^2 + 1)}{(x - 3)^2}
      =
      \frac{x^2 - 6x - 1}{(x - 3)^2}.
  \]

  \textbf{Verification.}
  The result is correct, with the final simplification added.

  \item \textbf{Stable item ID:} \texttt{ex-deriv-rules-004}.

  \textbf{Exercise.}
  Find the derivative of
  \[
    f(x) = \sqrt{3x^2 + 5x + 1}.
  \]

  \textbf{Solution.}
  By the chain rule,
  \[
    f'(x)
      =
      \frac{6x + 5}{2\sqrt{3x^2 + 5x + 1}}.
  \]

  \textbf{Verification.}
  The result is correct.

  \item \textbf{Stable item ID:} \texttt{ex-deriv-rules-005}.

  \textbf{Exercise.}
  Find the derivative of
  \[
    f(x) = \frac{x^2\sin x}{1+x^2}.
  \]

  \textbf{Solution.}
  Write \(f(x) = x^2 \cdot \frac{\sin x}{1+x^2}\). Then
  \[
    f'(x)
      =
      2x\frac{\sin x}{1+x^2}
      +
      x^2
      \frac{(1+x^2)\cos x - 2x\sin x}{(1+x^2)^2}.
  \]
  Combining over a common denominator gives
  \[
    f'(x)
      =
      \frac{2x\sin x + x^2(1+x^2)\cos x}{(1+x^2)^2}.
  \]

  \textbf{Verification.}
  The result is corrected and consolidated from the product/quotient split.

  \item \textbf{Stable item ID:} \texttt{ex-deriv-rules-006}.

  \textbf{Exercise.}
  Find the derivative of
  \[
    f(x) = e^{x^2}\cos x.
  \]

  \textbf{Solution.}
  By the product rule and chain rule,
  \[
    f'(x)
      =
      2xe^{x^2}\cos x - e^{x^2}\sin x
      =
      e^{x^2}(2x\cos x - \sin x).
  \]

  \textbf{Verification.}
  The result is correct.

  \item \textbf{Stable item ID:} \texttt{ex-deriv-rules-007}.

  \textbf{Exercise.}
  Find the derivative of
  \[
    f(x) = \ln(1+x^4).
  \]

  \textbf{Solution.}
  By the chain rule,
  \[
    f'(x)
      =
      \frac{4x^3}{1+x^4}.
  \]

  \textbf{Verification.}
  The result is correct.

  \item \textbf{Stable item ID:} \texttt{ex-deriv-rules-008}.

  \textbf{Exercise.}
  Find the derivative of
  \[
    f(x) = \frac{\sqrt{x^2+1}}{x^3 - 2x + 1}.
  \]

  \textbf{Solution.}
  Let
  \[
    g(x) = \sqrt{x^2+1},
    \qquad
    h(x) = x^3 - 2x + 1.
  \]
  Then
  \[
    g'(x) = \frac{x}{\sqrt{x^2+1}},
    \qquad
    h'(x) = 3x^2 - 2.
  \]
  By the quotient rule,
  \[
    f'(x)
      =
      \frac{
        (x^3 - 2x + 1)\frac{x}{\sqrt{x^2+1}}
        -
        \sqrt{x^2+1}(3x^2 - 2)
      }{(x^3 - 2x + 1)^2}.
  \]

  \textbf{Verification.}
  The result is corrected. The quotient-rule structure is right, and the
  radical derivative is \(x/\sqrt{x^2+1}\).
\end{enumerate}
